\documentclass{article}
\usepackage[utf8]{inputenc}
\usepackage{amsmath, amssymb}
\usepackage{graphicx}
\usepackage{xcolor}
\usepackage{geometry}
\usepackage{array}
\geometry{a4paper, margin=1in}

\begin{document}

\noindent
\textbf{Engenharia de Software / Engenharia de Requisitos}\\
\textbf{IFCE - Campus Boa Viagem}\\
\textbf{Trabalho de Engenharia de Software / Engenharia de Requisitos}\\
\textbf{Profa.:} Jéssica de Paulo Rodrigues (rodrigues.jessica@ifce.edu.br)\\
\textbf{Equipe:} Luidy Costa, Leticia Maciel, Germano Moraes\\
\textbf{Contato:} luidy.costa08@aluno.ifce.edu.br

\vspace{2cm}
\begin{center}
\textbf{Documentação do Sistema LungPrediction SaaS}\\[0.5cm]
\textbf{Especificação de Requisitos de Software}
\end{center}

\vspace{1cm}
\textbf{Histórico da Revisão}

\vspace{0.3cm}
\begin{tabular}{| m{3cm} | m{2cm} | m{7cm} | m{3cm} |}
\hline
\textbf{Data} & \textbf{Versão} & \textbf{Descrição} & \textbf{Autor} \\
\hline
23/Fev/26 & 1.0 & Criação do Documento ERS inicial & Luidy e Equipe \\
\hline
\end{tabular}

\newpage
\begin{center}
\textbf{Sumário}
\end{center}

\vspace{0.5cm}
\noindent
1 \textbf{Introdução} \dotfill 3\\
1.1 Objetivo \dotfill 3\\
1.2 Escopo \dotfill 3\\
1.3 Definições, Acrônimos e Abreviações \dotfill 3\\
1.4 Referências \dotfill 3\\
1.5 Visão Geral \dotfill 4\\
1.6 Perfil do usuário \dotfill 4\\

\vspace{0.3cm}
2 \textbf{Definição de Requisitos} \dotfill 4\\
2.1 Descrição do Propósito do Sistema \dotfill 4\\
2.2 Descrição do Minimundo \dotfill 4\\
2.3 Requisitos de Usuário \dotfill 4\\

\vspace{0.3cm}
3 \textbf{Identificação de Subsistemas}

\vspace{0.3cm}
4 \textbf{Modelo de Casos de Uso} \dotfill 6\\
5 \textbf{Protótipo} \dotfill 6

\newpage
\section{Introdução}
Este documento apresenta alguns dos resultados produzidos durante a fase de Especificação e Análise de Requisitos do sistema LungPrediction SaaS. A seção 2 aborda a documentação de requisitos, sendo apresentados o propósito do sistema (seção 2.1), a descrição do minimundo (seção 2.2) e a lista de requisitos de usuário levantados junto ao cliente (seção 2.3), definidos no formato de user stories, como requisitos funcionais e como requisitos não funcionais. A seção 3 apresenta os subsistemas identificados, mostrando suas dependências na forma de um diagrama de pacotes. A seção 4 apresenta o modelo de casos de uso, para representar uma visão comportamental do sistema. A seção 5 apresenta telas do protótipo construído para o sistema como forma de materializar os requisitos do usuário.

\subsection{Objetivo}
Especificar os requisitos do LungPrediction SaaS. A ERS descreve completamente o comportamento externo do aplicativo, que atua como uma plataforma de auxílio preditivo. Descreve também requisitos não funcionais, como segurança LGPD e arquitetura Multi-Tenant, necessários para fornecer uma descrição completa e abrangente dos requisitos para o software.

\subsection{Escopo}
O LungPrediction SaaS é uma aplicação Web operando no modelo Software as a Service (SaaS). O escopo central é fornecer uma triagem preditiva de risco de câncer de pulmão para pacientes, utilizando um modelo de Inteligência Artificial (Machine Learning) baseado em dados clínicos e sintomas. O sistema garante o isolamento de dados entre diferentes hospitais e a rastreabilidade dos laudos médicos através de geração de PDFs. Módulos de faturamento, agendamento de consultas ou controle de estoque hospitalar estão fora do escopo desta especificação.

\subsection{Definições, Acrônimos e Abreviações}
\begin{itemize}
    \item \textbf{SaaS}: Software as a Service (Software como Serviço).
    \item \textbf{Multi-Tenant}: Arquitetura onde uma única instância do software atende a múltiplos hospitais (inquilinos), mantendo os dados de cada um rigorosamente isolados.
    \item \textbf{LGPD}: Lei Geral de Proteção de Dados.
    \item \textbf{Upsert}: Lógica de banco de dados que atualiza um registro se ele existir, ou cria um novo se não existir.
\end{itemize}

\subsection{Referências}
\begin{itemize}
    \item IEEE Std 830-1998, IEEE Recommended Practice for Software Requirements Specifications.
    \item Conselho Federal de Medicina (CFM) - Resolução sobre Prontuários Eletrônicos e guarda de dados.
\end{itemize}
\section{Visão Geral}
O documento está organizado de forma a detalhar primeiramente o contexto do negócio clínico (Minimundo) e, em seguida, as especificações técnicas, histórias de usuário, arquitetura de subsistemas em MVC e os protótipos de tela.

\section{Perfil do usuário}
\begin{itemize}
    \item \textbf{Super Admin}: Desenvolvedores do sistema. Visualizam dados anonimizados para calibrar a IA.
    \item \textbf{Gestor do Hospital}: Conta jurídica (CNPJ) responsável por gerenciar a clínica e convidar médicos para a plataforma.
    \item \textbf{Médico}: usuário final (CRM). Interage com pacientes, insere sintomas e analisa o laudo da Inteligência Artificial.
\end{itemize}

\section{Definição de Requisitos}
\subsection{Descrição do Propósito do Sistema}
O propósito geral do sistema é atuar como uma segunda opinião analítica de alta velocidade, processando sinais vitais e anamneses de pacientes para calcular, matematicamente, a probabilidade de risco para o câncer de pulmão, auxiliando médicos na tomada de decisão rápida em ambientes hospitalares.

\subsection{Descrição do Minimundo}
Em ambientes clínicos e hospitalares, a triagem de pacientes com sintomas respiratórios é um processo crítico. Médicos coletam dados como ocorrência de falta de ar, tosse com sangue, hábitos de fumo e sinais vitais. Muitas vezes, correlações sutis entre múltiplos fatores podem passar despercebidas na correria do plantão, atrasando diagnósticos de câncer de pulmão. O sistema resolve esse problema fornecendo uma interface rápida onde o médico insere o prontuário. O modelo de Machine Learning traduz esses dados e retorna o risco. Como os médicos costumam dar plantão em várias unidades diferentes, o sistema precisa permitir que um único profissional (identificado pelo seu CRM) se vincule a vários hospitais, garantindo que os dados de pacientes de um hospital jamais sejam vistos por outro (Multi-Tenant). Além disso, por exigências legais, prontuários gerados não podem ser deletados fisicamente do banco de dados.

\subsection{Requisitos de Usuário}
Os requisitos de usuário identificados para o sistema são apresentados a seguir, no formato de \textit{user stories}:

\begin{tabular}{|l|l|l|p{6cm}|p{6cm}|}
\hline
\textbf{ID} & \textbf{Depende de} & \textbf{Prioridade} & \textbf{Descrição} & \textbf{Critérios de Aceitação} \\
\hline
US01 & Nenhuma & Alta & Eu como Gestor quero cadastrar médicos para formar minha equipe. & 1.1 O sistema deve disparar e-mail de convite. 1.2 O vínculo é ativado após aceite. \\
\hline
US02 & US01 & Alta & Eu como Médico quero selecionar um hospital para ver os pacientes do plantão. & 2.1 O painel deve listar apenas prontuários da unidade selecionada. \\
\hline
US03 & US02 & Alta & Eu como Médico quero inserir sintomas para obter o risco de câncer. & 3.1 A IA deve retornar a predição. 3.2 Risco $\geq$ 50\% é ``Alto Risco''. \\
\hline
US04 & US03 & Média & Eu como Médico quero exportar a triagem para entregar laudo formal. & 4.1 O PDF deve conter CRM, sintomas marcados e a \% de risco. \\
\hline
\end{tabular}

\vspace{0.5cm}

\begin{tabular}{|l|p{8cm}|l|l|}
\hline
\textbf{ID} & \textbf{Descrição} & \textbf{Categoria} & \textbf{Req. Relacionados} \\
\hline
RF01 & Processar dados clínicos (Dicionário) e traduzir em matriz binária para consumo da IA. & Essencial & US03 \\
\hline
RF02 & Utilizar lógica de Upsert: atualizar histórico se existir ou criar novo se ID vazio. & Essencial & US03 \\
\hline
RF03 & Gerar documentos PDF estáticos com o resumo da triagem. & Importante & US04 \\
\hline
\end{tabular}

\end{document}